\documentclass[11pt,letterpaper]{article}
\usepackage[margin=0.65in]{geometry}
\usepackage[T1]{fontenc}
\usepackage{lmodern,amsmath,amssymb,microtype}
\usepackage{titlesec}
\titleformat{\section}{\normalfont\normalsize\bfseries}{}{0pt}{}
\titlespacing*{\section}{0pt}{6pt}{3pt}
\usepackage[colorlinks=true,urlcolor=blue,citecolor=blue,linkcolor=blue]{hyperref}
\hypersetup{pdftitle={Updating Nystrom Preconditioners with Solver Residuals}}
\pagestyle{empty}
\setlength{\parindent}{0pt}
\setlength{\parskip}{4pt}
\newcommand{\nystrom}{Nystr\"om}
\begin{document}
\fontsize{10.5}{12.3}\selectfont
\begin{center}
  {\large\bfseries Updating Nystr\"om Preconditioners with Solver Residuals}\\[3pt]
  CS 6220: Data-sparse Matrix Computations \quad Project Proposal
\end{center}
\vspace{-7pt}

\textbf{Research question and motivation.}
Kernel ridge regression learns a predictor by solving linear equations built from
pairwise similarities between training examples.
A \nystrom{} preconditioner approximates this similarity matrix using a small set
of its columns and accelerates the conjugate-gradient solver~\cite{frangella2023nystrom}.
Keeping more columns costs more memory; with a limited column budget, the choice
of columns matters. A set chosen to approximate the whole matrix may spend that
budget on components of the error that the solver has already corrected.
The residual, the mismatch between the two sides of the equations, provides
information about what remains to be solved.
\textbf{Can the current residual guide kernel-column replacements that reduce the
total time to reach a prescribed solution accuracy, while keeping the number of
retained columns fixed?}
The hypothesis is that replacing columns as the residual changes can make a
small preconditioner more effective throughout the solve.

Task-dependent kernel approximation already uses prediction targets to guide
column selection~\cite{bach2005predictive}, and limited-memory preconditioning
uses information from iterative solves~\cite{morales2000automatic}.
I will develop a rule that uses the changing residual to decide which kernel
columns to discard and which to acquire. Any reduction in iteration time must
outweigh the cost of column replacement and solver restarts.

\textbf{Initial method.}
Let $A=K+\lambda I$, where $K$ is a Gaussian kernel matrix and $\lambda>0$ is the
regularization parameter. The system $A\alpha=y$ determines the prediction
coefficients $\alpha$ from training targets $y$.
The initial $k$ columns will be selected by randomly pivoted
Cholesky~\cite{chen2025rpcholesky}. For their indices $S$, the preconditioner is
\[
 M_S=K_{:,S}K_{S,S}^{\dagger}K_{S,:}+\lambda I,
\]
where $\dagger$ denotes the pseudoinverse.
Here $K_{:,S}$ contains the selected columns and $K_{S,S}$ their entries at the
selected rows. Retaining $k$ columns requires $O(nk)$ storage for $n$ examples.

After a block of conjugate-gradient iterations, the method computes the residual
$r_t=y-A\alpha_t$ and samples a small number of candidate column indices with
probabilities proportional to $r_{t,i}^2$, emphasizing equations with larger
remaining discrepancies.
For each candidate, it considers replacing one retained column and selects the set
$S'$ with the smallest score $r_t^{\mathsf T}M_{S'}^{-1}r_t$.
The score is an upper bound on $r_t^{\mathsf T}A^{-1}r_t=e_t^{\mathsf T}Ae_t$,
where $e_t$ is the difference between the exact and current solutions.
A smaller score tightens this estimate of the remaining solution error.
The experiment tests whether the replacements it favors improve subsequent convergence.
An exchange is accepted when it reduces the score, keeping $k$ retained columns.
The solver then restarts from the current solution.
Updates of matrices of size $k+1$ allow exchanges to be evaluated using the
retained columns and one candidate column at a time.
I will revise the method as needed based on the convergence and timing results.

\textbf{Experiments and interpretation.}
I will evaluate the method on clustered synthetic inputs and the UCI
\href{https://archive.ics.uci.edu/dataset/165/concrete+compressive+strength}{Concrete}
and \href{https://archive.ics.uci.edu/dataset/291/airfoil+self+noise}{Airfoil}
regression datasets, with at most 2,048 training examples.
The experiments will vary the column budget, regularization, and replacement frequency.
Changing the synthetic targets while keeping the kernel fixed will test
whether replacement adapts to the particular system being solved.
Comparisons will include fixed and periodically reconstructed \nystrom{} preconditioners,
and limited-memory quasi-Newton preconditioning with comparable vector storage.
Exchanges based on the initial targets instead of the current residual test the
value of the current residual; restarting without exchanges isolates restart costs.

All solvers will stop at a directly checked relative residual of $10^{-8}$.
Measurements will include total runtime on CPUs, replacement overhead,
matrix--vector products, and peak memory over five seeds.
Direct Cholesky solves will check solution and prediction accuracy.
Residual histories and score changes will help assess whether accepted
replacements improve subsequent convergence.

\vspace{-2pt}
{\footnotesize
\setlength{\parskip}{0pt}
\bibliographystyle{unsrt}
\bibliography{references}
}
\vspace{-3pt}
{\scriptsize AI assistance: OpenAI Codex assisted with literature search and proposal drafting.}
\end{document}
